% !PRINTED 370
\weberpage{370}

The difference (5) sinks below every limit as $n$ increases
without bound, and since $x_n > a_n$, and hence $Q_n x_n + Q_{n-1} > Q_{n+1}$,
it is in absolute value smaller than
\weq{6}{\frac{1}{Q_n Q_{n+1}}}

and all the more smaller than $1 : Q_n^2$.

The numbers of the sequence (3) thus approach the limit $x$
decreasingly, those of the sequence (4) increasingly. The expression (6)
gives a measure of the error one commits when $x$ is replaced by
the convergent $P_n : Q_n$.

The convergents are thus approximate expressions
of irrational numbers by rational
fractions.

That these convergents, for an equal degree of approximation
to the irrational number $x$, are the simplest possible
is expressed by the following theorem:

\gesperrt{IV. Between the two rational
fractions}
\[
\frac{P_n}{Q_n}, \quad \frac{P_{n-1}}{Q_{n-1}}
\]
\gesperrt{no other rational fraction can be inserted
whose denominator is smaller than $Q_n$ or even only
equal to $Q_n$.}

Suppose that the rational fraction $M : N$ lies between
the convergents $P_n : Q_n$ and $P_{n-1} : Q_{n-1}$. Then the
difference
\[
\frac{P_n}{Q_n} - \frac{P_{n-1}}{Q_{n-1}}
\]
is absolutely greater than, and of the same sign as,
\[
\frac{M}{N} - \frac{P_{n-1}}{Q_{n-1}},
\]
hence
\[
(-1)^n \left(\frac{P_n}{Q_n} - \frac{P_{n-1}}{Q_{n-1}}\right) > (-1)^n \left(\frac{M}{N} - \frac{P_{n-1}}{Q_{n-1}}\right).
\]

If one multiplies both sides by $N Q_{n-1}$, then by (1) it follows that
\[
\frac{N}{Q_n} > (-1)^n (M Q_{n-1} - N P_{n-1}),
\]
and since on the right stands a positive integer, which is therefore at least
equal to 1, it follows that
\[
N > Q_n.
\]
